% Lucas's completed solution, reviewed on October 8, 2026, at his request.
% Red marks revised prose, mathematical corrections, and added justifications.
% Purely typographical changes to math delimiters are not marked.
% All five supplied parts have now been reviewed; retain the red review marks.
\noindent\textbf{Solution.} For part (a), \textcolor{red}{suppose} a matrix
\(P_\sigma\) is a permutation matrix where \(c_k=e_{\sigma(k)}\).
Since \(\sigma:\{1,\ldots,n\}\to\{1,\ldots,n\}\) is bijective,
\(\sigma(1),\ldots,\sigma(n)\) will be distinct and
\textcolor{red}{exhaust} \(\{1,\ldots,n\}\).
\textcolor{red}{Each column is a standard basis vector, so it contains a
single \(1\) in row \(\sigma(k)\) and zeros elsewhere. Thus every column
contains exactly one \(1\). Since each row index occurs exactly once among
\(\sigma(1),\ldots,\sigma(n)\), every row also contains exactly one \(1\).}

Conversely, consider \textcolor{red}{an \(n\times n\) matrix} \(P\)
where every entry is \(1\) or \(0\) and there is exactly one \(1\)
in each row and column. For each column \(k\), there exists
\textcolor{red}{a unique index} \(\sigma(k)\)
\textcolor{red}{giving the row containing its \(1\). Hence}
\[
  c_k=e_{\sigma(k)}.
\]

No two columns can have their \(1\) in the same row, so \(\sigma\) is
injective. \textcolor{red}{An injective map from a finite set to itself
is also surjective, so \(\sigma\) is bijective.} Thus \(\sigma\in S_n\)
and
\[
  P=\begin{bmatrix}e_{\sigma(1)}&\ldots&e_{\sigma(n)}\end{bmatrix}
    =P_\sigma.
\]
Therefore \(P\) is a permutation matrix if and only if every entry is
\(0\) or \(1\) with exactly one \(1\) in each row and column. \qed

For part (b), we know that when we multiply a matrix \(M\) by a basis
vector it produces \textcolor{red}{the corresponding column},
\(Me_i=c_i\). \textcolor{red}{Matrix multiplication acts separately on
each column of the matrix on the right.} So when multiplying \(M\) by
the matrix of basis vector columns,
\[
  M\begin{bmatrix}e_1\mid\ldots\mid e_n\end{bmatrix}
    =\begin{bmatrix}c_1\mid\ldots\mid c_n\end{bmatrix}.
\]
\textcolor{red}{Applying the same rule to the permuted basis columns gives}
\begin{align*}
  M\begin{bmatrix}e_{\sigma(1)}\mid\ldots\mid e_{\sigma(n)}\end{bmatrix}
    &=\textcolor{red}{\begin{bmatrix}Me_{\sigma(1)}\mid\ldots\mid
      Me_{\sigma(n)}\end{bmatrix}}\\
    &=\begin{bmatrix}c_{\sigma(1)}\mid\ldots\mid c_{\sigma(n)}\end{bmatrix}.
\end{align*}
Since \(P_\sigma=\begin{bmatrix}e_{\sigma(1)}\mid\ldots\mid
e_{\sigma(n)}\end{bmatrix}\),
\[
  MP_\sigma=\begin{bmatrix}c_{\sigma(1)}\mid\ldots\mid
    c_{\sigma(n)}\end{bmatrix}.
\]
\qed

For part (c), we first express \(v\in\R^n\) as a linear combination
of \textcolor{red}{the standard basis vectors},
\[
  v=v^1e_1+\ldots+v^ne_n.
\]
Now applying \(P_\sigma\) to both sides and using linearity,
\begin{align*}
  P_\sigma v&=P_\sigma(v^1e_1+\ldots+v^ne_n)\\
    &=v^1P_\sigma e_1+\ldots+v^nP_\sigma e_n.
\end{align*}
Focusing on each \(k\)-th term \(P_\sigma e_k\), we know that
right-multiplying by the \(k\)-th basis vector produces the \(k\)-th
column. \textcolor{red}{By the definition of \(P_\sigma\), this gives}
\[
  P_\sigma e_k=\textcolor{red}{e_{\sigma(k)}}.
\]
\textcolor{red}{Substituting this identity into the linear combination gives}
\begin{align*}
  \textcolor{red}{P_\sigma v}
    &=v^1e_{\sigma(1)}+\ldots+v^n\textcolor{red}{e_{\sigma(n)}}\\
    &=\sum_{k=1}^n v^ke_{\sigma(k)}.
\end{align*}
Thus, \(P_\sigma\) defines the linear map that permutes the coordinates
of \(v\) according to \(\sigma\), \textcolor{red}{moving} the coefficient
\(v^k\) from the \(k\)-th coordinate to the \(\sigma(k)\)-th coordinate.
\textcolor{red}{Equivalently, the \(i\)-th output coordinate comes from
the input index \(k\) satisfying \(\sigma(k)=i\), so}
\[
  \textcolor{red}{(P_\sigma v)^i=v^{\sigma^{-1}(i)}.}
\]
\qed

For part (d), \textcolor{red}{fix any \(k\in\{1,\ldots,n\}\).}
It was proven in part (c) that
\[
  P_\tau e_k=e_{\tau(k)}.
\]
\textcolor{red}{Applying \(P_\sigma\) and using its action on standard
basis vectors gives}
\begin{align*}
  P_\sigma P_\tau e_k&=P_\sigma e_{\tau(k)}\\
    &=e_{\sigma(\tau(k))}.
\end{align*}
\textcolor{red}{On the other hand, the composition satisfies
\((\sigma\circ\tau)(k)=\sigma(\tau(k))\), so}
\begin{align*}
  P_{\sigma\circ\tau}e_k
    &=\textcolor{red}{e_{(\sigma\circ\tau)(k)}}\\
    &=e_{\sigma(\tau(k))}.
\end{align*}
\textcolor{red}{Thus \(P_{\sigma\circ\tau}e_k=P_\sigma P_\tau e_k\)
for every \(k\). Every vector is a linear combination of the standard
basis vectors, so linearity implies that the two matrices act identically
on every vector. Therefore,}
\[
  P_{\sigma\circ\tau}=P_\sigma P_\tau.
\]
\qed

% Keep the final determinant argument together in the handout and appendix.
\newpage
For part (e), the determinant of \textcolor{red}{\(P_\sigma\)} is
\textcolor{red}{the value of the normalized antisymmetric multilinear
function \(D\) on its columns. Here \(c_k=e_{\sigma(k)}\), so}
\begin{align*}
  \textcolor{red}{\det(P_\sigma)}&=D(c_1,\ldots,c_n)\\
    &=D(e_{\sigma(1)},\ldots,e_{\sigma(n)}).
\end{align*}
\textcolor{red}{Each column is already a standard basis vector, so there
is just one basis-vector value to evaluate. Since \(\sigma\) is bijective,
these inputs are distinct.}

\textcolor{red}{Choose a decomposition
\(\sigma=\tau_1\circ\cdots\circ\tau_r\) into \(r\) transpositions.}
Whenever two \textcolor{red}{inputs are interchanged in an antisymmetric
multilinear function}, a factor of \(-1\) is produced.
\textcolor{red}{Applying the transpositions in this decomposition
rearranges \((e_1,\ldots,e_n)\) into
\((e_{\sigma(1)},\ldots,e_{\sigma(n)})\), contributing a factor
of \((-1)^r\). Normalization gives \(D(e_1,\ldots,e_n)=1\), and
the definition of the sign gives \(\epsilon(\sigma)=(-1)^r\). Hence,}
\begin{align*}
  D(e_{\sigma(1)},\ldots,e_{\sigma(n)})
    &=\textcolor{red}{(-1)^r}D(e_1,\ldots,e_n)\\
    &=\textcolor{red}{(-1)^r\cdot1}\\
    &=\epsilon(\sigma).
\end{align*}
\textcolor{red}{Therefore, \(\det(P_\sigma)=\epsilon(\sigma)\).} \qed
